% Options for packages loaded elsewhere
% Options for packages loaded elsewhere
\PassOptionsToPackage{unicode}{hyperref}
\PassOptionsToPackage{hyphens}{url}
\PassOptionsToPackage{dvipsnames,svgnames,x11names}{xcolor}
%
\documentclass[
  letterpaper,
  DIV=11,
  numbers=noendperiod]{scrartcl}
\usepackage{xcolor}
\usepackage{amsmath,amssymb}
\setcounter{secnumdepth}{5}
\usepackage{iftex}
\ifPDFTeX
  \usepackage[T1]{fontenc}
  \usepackage[utf8]{inputenc}
  \usepackage{textcomp} % provide euro and other symbols
\else % if luatex or xetex
  \usepackage{unicode-math} % this also loads fontspec
  \defaultfontfeatures{Scale=MatchLowercase}
  \defaultfontfeatures[\rmfamily]{Ligatures=TeX,Scale=1}
\fi
\usepackage{lmodern}
\ifPDFTeX\else
  % xetex/luatex font selection
\fi
% Use upquote if available, for straight quotes in verbatim environments
\IfFileExists{upquote.sty}{\usepackage{upquote}}{}
\IfFileExists{microtype.sty}{% use microtype if available
  \usepackage[]{microtype}
  \UseMicrotypeSet[protrusion]{basicmath} % disable protrusion for tt fonts
}{}
\makeatletter
\@ifundefined{KOMAClassName}{% if non-KOMA class
  \IfFileExists{parskip.sty}{%
    \usepackage{parskip}
  }{% else
    \setlength{\parindent}{0pt}
    \setlength{\parskip}{6pt plus 2pt minus 1pt}}
}{% if KOMA class
  \KOMAoptions{parskip=half}}
\makeatother
% Make \paragraph and \subparagraph free-standing
\makeatletter
\ifx\paragraph\undefined\else
  \let\oldparagraph\paragraph
  \renewcommand{\paragraph}{
    \@ifstar
      \xxxParagraphStar
      \xxxParagraphNoStar
  }
  \newcommand{\xxxParagraphStar}[1]{\oldparagraph*{#1}\mbox{}}
  \newcommand{\xxxParagraphNoStar}[1]{\oldparagraph{#1}\mbox{}}
\fi
\ifx\subparagraph\undefined\else
  \let\oldsubparagraph\subparagraph
  \renewcommand{\subparagraph}{
    \@ifstar
      \xxxSubParagraphStar
      \xxxSubParagraphNoStar
  }
  \newcommand{\xxxSubParagraphStar}[1]{\oldsubparagraph*{#1}\mbox{}}
  \newcommand{\xxxSubParagraphNoStar}[1]{\oldsubparagraph{#1}\mbox{}}
\fi
\makeatother


\usepackage{longtable,booktabs,array}
\newcounter{none} % for unnumbered tables
\usepackage{calc} % for calculating minipage widths
% Correct order of tables after \paragraph or \subparagraph
\usepackage{etoolbox}
\makeatletter
\patchcmd\longtable{\par}{\if@noskipsec\mbox{}\fi\par}{}{}
\makeatother
% Allow footnotes in longtable head/foot
\IfFileExists{footnotehyper.sty}{\usepackage{footnotehyper}}{\usepackage{footnote}}
\makesavenoteenv{longtable}
\usepackage{graphicx}
\makeatletter
\newsavebox\pandoc@box
\newcommand*\pandocbounded[1]{% scales image to fit in text height/width
  \sbox\pandoc@box{#1}%
  \Gscale@div\@tempa{\textheight}{\dimexpr\ht\pandoc@box+\dp\pandoc@box\relax}%
  \Gscale@div\@tempb{\linewidth}{\wd\pandoc@box}%
  \ifdim\@tempb\p@<\@tempa\p@\let\@tempa\@tempb\fi% select the smaller of both
  \ifdim\@tempa\p@<\p@\scalebox{\@tempa}{\usebox\pandoc@box}%
  \else\usebox{\pandoc@box}%
  \fi%
}
% Set default figure placement to htbp
\def\fps@figure{htbp}
\makeatother


% definitions for citeproc citations
\NewDocumentCommand\citeproctext{}{}
\NewDocumentCommand\citeproc{mm}{%
  \begingroup\def\citeproctext{#2}\cite{#1}\endgroup}
\makeatletter
 % allow citations to break across lines
 \let\@cite@ofmt\@firstofone
 % avoid brackets around text for \cite:
 \def\@biblabel#1{}
 \def\@cite#1#2{{#1\if@tempswa , #2\fi}}
\makeatother
\newlength{\cslhangindent}
\setlength{\cslhangindent}{1.5em}
\newlength{\csllabelwidth}
\setlength{\csllabelwidth}{3em}
\newenvironment{CSLReferences}[2] % #1 hanging-indent, #2 entry-spacing
 {\begin{list}{}{%
  \setlength{\itemindent}{0pt}
  \setlength{\leftmargin}{0pt}
  \setlength{\parsep}{0pt}
  % turn on hanging indent if param 1 is 1
  \ifodd #1
   \setlength{\leftmargin}{\cslhangindent}
   \setlength{\itemindent}{-1\cslhangindent}
  \fi
  % set entry spacing
  \setlength{\itemsep}{#2\baselineskip}}}
 {\end{list}}
\usepackage{calc}
\newcommand{\CSLBlock}[1]{\hfill\break\parbox[t]{\linewidth}{\strut\ignorespaces#1\strut}}
\newcommand{\CSLLeftMargin}[1]{\parbox[t]{\csllabelwidth}{\strut#1\strut}}
\newcommand{\CSLRightInline}[1]{\parbox[t]{\linewidth - \csllabelwidth}{\strut#1\strut}}
\newcommand{\CSLIndent}[1]{\hspace{\cslhangindent}#1}



\setlength{\emergencystretch}{3em} % prevent overfull lines

\providecommand{\tightlist}{%
  \setlength{\itemsep}{0pt}\setlength{\parskip}{0pt}}



 


\KOMAoption{captions}{tableheading}
\makeatletter
\makeatother
\makeatletter
\@ifpackageloaded{caption}{}{\usepackage{caption}}
\@ifpackageloaded{subcaption}{}{\usepackage{subcaption}}
\makeatother
\usepackage{bookmark}
\IfFileExists{xurl.sty}{\usepackage{xurl}}{} % add URL line breaks if available
\urlstyle{same}
% fallback for those not using the hyperref driver hyperxmp:
\makeatletter
\@ifundefined{xmpquote}{\newcommand{\xmpquote}[1]{#1}}{}
\makeatother
\hypersetup{
  pdftitle={Visualizing Educational Attainment as a Lever for Upward Mobility in the District of Columbia},
  pdfauthor={Úna Gogstetter; Ayla Hauser; Simran Jayasinghe; Alice MacGarvie Thompson},
  pdfkeywords={\xmpquote{Washington, D.C}, \xmpquote{Upward
Mobility}, \xmpquote{Educational Attainment}, \xmpquote{Data
Dashboard}},
  colorlinks=true,
  linkcolor={blue},
  filecolor={Maroon},
  citecolor={Blue},
  urlcolor={Blue},
  pdfcreator={LaTeX via pandoc}}


\title{Visualizing Educational Attainment as a Lever for Upward Mobility
in the District of Columbia}
\author{Úna Gogstetter \and Ayla Hauser \and Simran
Jayasinghe \and Alice MacGarvie Thompson}
\date{2026-11-08}
\begin{document}
\maketitle
\begin{abstract}
The drastic economic disparities in Washington, District of Columbia,
often obscure the impacts of poverty and inequity upon residents living
in working-class and low-income communities. The immense wealth in parts
of the city skews citywide data and masks these impacts. This limits the
capability of city officials to use data to inform education policy.

Many data-driven projects do not examine data with such granularity as
the school, census tract, block, or neighborhood level. Projects such as
the Urban Institute's Upward Mobility Dashboard help policymakers
understand conditions affecting upward mobility. However, many of these
projects only cover data at the county level, not the census tract
level.

Our project seeks to visualize and model the impacts of key educational
conditions and levers on upward mobility and educational attainment
outcomes. We use data pulled from the school, block, tract,
neighborhood, and ward levels to model this for the Department of the
Deputy Mayor for Education for the District of Columbia.
\end{abstract}


\section{Introduction}\label{introduction}

\begin{quote}
Introduction and Motivation: What is the problem you are trying to
solve? Why is this problem interesting? What has been tried before? What
have been the shortcoming of those approaches that necessitate your
efforts? Often, this section will conclude with a subsection (or
paragraph) outlining ``our contributions.'' What is the new knowledge
that this paper contributes?
\end{quote}

In the District of Columbia, significant disparities in resources and
economic mobility outcomes create different realities of opportunity for
residents across the district. Aggregate citywide statistics often mask
contrasts between local metrics, losing track of inequities through
broad averages of the data, an oversimplification that can be observed
when exploring poverty and income across the district.

D.C. exhibits one of the highest levels of income inequality in the
United States, where household income in the highest quintile exceeds
that of the lowest quintile by nearly 29 times (Santillo et al. 2022).
D.C.'s Gini coefficient --- which is a value between 0 and 1 where 1
signifies the greatest income inequality --- is 0.51 (U.S. Census Bureau
2023). The county encompassing the District of Columbia is above the
95th percentile for income inequality (as measured by the Gini
coefficient) among all counties in the United States (U.S. Census Bureau
2023). These economic gaps are heavily spatialized across the city's
eight wards; for instance, poverty concentrations predominantly cluster
east of the Anacostia River in Wards 7 and 8 (Naveed 2017). Since
commonly reported citywide metrics average out these regional divides,
effective policy design requires granular spatial analysis to reveal how
opportunity drivers and community resources are distributed across
specific neighborhoods.

Because educational access and environment serve as catalysts for upward
economic mobility, these systemic resource disparities directly impact
outcomes across D.C.'s public education system. To address these
structural inequities, the Office of the District of Columbia Deputy
Mayor for Education (DME) works to create conditions that support
academic acceleration across the PreK--12 public education system. The
Office of the District of Columbia Deputy Mayor for Education (DME)
works to create conditions that support educational attainment and
academic acceleration at the PreK-12 levels, focusing on schools and
their surrounding neighborhoods as sites for social and economic
justice. This work relies on granular place-based insights to inform
targeted policymaking, funding initiatives, and community programming,
and ensure these programs are guided to areas of need.

National data visualization platforms such as the Urban Institute's
Upward Mobility Data Dashboard provide helpful insights into key
predictors of mobility indicators, but lack the neighborhood resolution
required for D.C.'s urban landscape. Without multi-scalar analytical
tools that synthesize local education and socio-economic outcomes
simultaneously, government agencies face challenges in identifying where
and how targeted interventions can close persistent opportunity gaps.

To bridge this gap, this capstone project collaborates with the DME to
adapt and refine the Urban Institute's Upward Mobility Dashboard at the
local level. This project specifically explores the impact of
educational factors on upward mobility across the District of Columbia
by aggregating disparate indicators into a unified structure across
census tract, neighborhood, ward, and citywide scales. The project aims
to create a public-facing web application that illustrates variations in
opportunity and offers statistical analysis tools that can be used to
evaluate how targeted policy interventions can close gaps in
opportunity, providing a localized platform for public reporting,
equitable resource allocation, and data-driven decision-making.

\subsection{Relevant work}\label{relevant-work}

\subsection{Our contribution}\label{our-contribution}

\begin{verbatim}
[1] "C:/Users/unael/Documents/dced-f26/quarto-manuscript"
\end{verbatim}

\textsubscript{Source:
\href{https://sds-capstone.github.io/dced-f26/test.qmd.html}{Article
Notebook}}

\protect\phantomsection\label{refs}
\begin{CSLReferences}{1}{1}
\bibitem[\citeproctext]{ref-naveed2017income}
Naveed, Minahil. 2017. {``Income Inequality in {DC} Highest in the
Country.''} DC Fiscal Policy Institute, December 15.
\url{https://dcfpi.org/all/income-inequality-dc-highest-country-2/}.

\bibitem[\citeproctext]{ref-santillo2022washingtondc}
Santillo, Miranda, Mingli Zhong, and Oriya Cohen. 2022. {``Washington,
{DC}, Has Glaring Financial Health Inequities. So Do Most {American}
Cities. Urban Wire.''} Urban Institute, October 6.
\url{https://www.urban.org/urban-wire/washington-dc-has-glaring-financial-health-inequities-so-do-most-american-cities}.

\bibitem[\citeproctext]{ref-Census2023ACSDT5Y2023}
U.S. Census Bureau. 2023. {``Gini Index of Income Inequality.''} U.S.
Department of Commerce.
\url{https://data.census.gov/table/ACSDT5Y2023.B19083?q=income+inequality&g=010XX00US\%240500000}.

\end{CSLReferences}




\end{document}
